% ===== Block B16: 原文 p.151–160 | §6.5 $A_n^{(1)}$ 与 Hecke 代数 | 译手 tr-ch6b =====
% 说明：§6.5 的标题起自原文 p.151 中部，该页前半为 §6.4 的末尾（属 B15），本块不译。
% [原文 p.151]
\subsection{$A_n^{(1)}$ 与 Hecke 代数}

以下我们讨论复单代数 $\mathfrak g$ 为 $A_n^{(1)}$ 型的情形，即
$\mathfrak g=\mathfrak{sl}(n+1,\C)$。

\noindent\textbf{6.5.1 杨图.}\quad
首先考虑李代数 $\mathfrak{su}(n+1)$（或 $\mathfrak{sl}(n+1,\C)$）的支配整权
（dominant integral weight）与杨图（Young diagram）之间的关系。
$\mathfrak{su}(n+1)$ 的元素由 $(n+1)\times(n+1)$ 无迹（traceless）矩阵表示。设
$\varepsilon_i$ 为除第 $(i,i)$ 位置为 $1$ 之外其余位置全为 $0$ 的矩阵。于是
$\mathfrak{sl}(n+1,\C)$ 的一组单根（simple root）由
\[
\alpha_i=\varepsilon_i-\varepsilon_{i+1},\quad i=1,2,\dots,n
\]
给出。根格（root lattice）描述为
\[
Q=\sum_{i=1}^{n+1} n_i\varepsilon_i,\quad n_i\in\Z,\quad
\sum_{i=1}^{n+1}n_i=0.
\]
归一化 Killing 型 $(\ ,\ )$ 是由内积（inner product）
\begin{equation}
\varepsilon_i\cdot\varepsilon_j=\delta_{ij}
\tag{6.34}
\end{equation}
诱导得到的。

% [Tutor Note: $\varepsilon_i$ 是第 $(i,i)$ 位置为 $1$ 的矩阵，它不是反 Hermite 的，严格说来
% 属于 $\mathfrak{gl}(n+1)$ 而不属于 $\mathfrak{su}(n+1)$；原文是说 $\mathfrak{su}(n+1)$
% 的复化（即 $\mathfrak{sl}(n+1,\C)$）中的元素用无迹矩阵表示，$\varepsilon_i$ 只在
% $\mathfrak{sl}(n+1,\C)$ 中作为对角矩阵使用。$(\ ,\ )$ 在 $\mathfrak{sl}(n+1,\C)$ 上
% 即 $\Tr(XY)$ 按 $n+1$ 归一化后的形式，故 $\varepsilon_i\cdot\varepsilon_j=\delta_{ij}$。

% [原文 p.152]
\figomit{6.2}{杨图 $[5,4,3,1,1]$。}

基本权（fundamental weight）$\omega_i$，$i=1,2,\dots,n+1$，由
\[
\omega_i=\varepsilon_1+\varepsilon_2+\cdots+\varepsilon_i
-\frac{i}{n+1}\sum_{k=1}^{n+1}\varepsilon_k
\]
给出。于是，支配整权 $\lambda$ 可写为
\begin{equation}
\lambda=\sum_{i=1}^n\lambda_i\omega_i
=\sum_{j=1}^{n+1}(f_j-\kappa(\lambda))\varepsilon_j,
\quad \lambda_i\in\Z_{\ge0}
\tag{6.35}
\end{equation}
其中
\begin{align}
f_j &= \lambda_i+\lambda_{i+1}+\cdots+\lambda_{n+1},
\quad j=1,\dots,n, \tag{6.36}\\
f_{n+1} &= 0, \nonumber\\
\kappa(\lambda) &= \frac{1}{n+1}\sum_{j=1}^n j\lambda_j. \nonumber
\end{align}
% [Tutor Note: 式 (6.36) 右端首项原文印作 $\lambda_i$，按定义应为 $\lambda_j$（$f_j$ 是第
% $j$ 行之格数）；末项 $\lambda_{n+1}$ 亦然，因 $\lambda=\sum_{i=1}^n\lambda_i\omega_i$
% 只有 $n$ 个分量、且约定 $f_{n+1}=0$，此处应理解为求和到 $\lambda_n$。以上照原文排印。

现在，对一个支配整权 $\lambda=\sum_{i=1}^n\lambda_i\omega_i$，我们可以联系一个杨图
$[f_1,f_2,\dots,f_n]$：它有 $n$ 行，第 $j$ 行由 $f_j$ 个格子组成，其中 $f_j$ 由上述
公式 (6.36) 给出。注意最高根（highest root）$\theta$ 由 $\varepsilon_1-\varepsilon_{n+1}$
给出，于是由 (6.34) 与 (6.35)，
\[
(\lambda,\theta)\le\ell
\]
等价于
\[
f_1\le\ell.
\]
因此杨图的集合
\[
\Gamma_{n,\ell}=\{[f_1,\dots,f_n]\mid f_1\le\ell\}
\]
可以看作行数不超过 $n$、且对应权满足 $(\lambda,\theta)\le\ell$ 的一切杨图的集合。
于是 $\Gamma_{n,\ell}$ 与标号集（labelling set）$P_\ell$ 之间存在一一对应。

更一般地，对深度（depth）不超过 $n+1$ 的杨图 $Y=[f_1,f_2,\dots,f_{n+1}]$，置
\[
\begin{aligned}
\lambda(Y)&=\sum_{j=1}^n b_j\omega_j,\qquad b_j=f_j-f_{j+1},\\
\ell(Y)&=f_1-f_{n+1}=\sum_{j=1}^n b_j=(\lambda(Y),\theta),\\
m(Y)&=\frac{1}{n+1}\sum_{j=1}^{n+1}f_j
=\frac{|Y|}{n+1}+f_{n+1}+\frac{1}{n+1}\sum_{j=k} kb_k.
\end{aligned}
\]
% [Tutor Note: $m(Y)$ 一行末项的求和号原文印作 $\sum_{j=k}$，下标残缺；由
% $\sum_{j=1}^{n+1}f_j=\sum_{k=1}^n kb_k+(n+1)f_{n+1}$ 可知应为 $\sum_{k=1}^{n}kb_k$。
% 上式照原文排印，读者按此理解。$|Y|=\sum_{j=1}^{n+1}f_j$ 是杨图的总格数。

于是对 $\lambda\in P_\ell$，存在深度 $\le n+1$ 的杨图 $Y$，使 $\ell(Y)=\ell$ 且
$\lambda(Y)=\lambda$。杨图 $Y$ 不唯一，因为
$[\underbrace{1,1,\dots,1}_{n+1}]$ 对应平凡表示。注意对每个 $i$，数 $f_i-m(Y)$
与满足 $\lambda(Y)=\lambda$ 的杨图 $Y$ 的选取无关，因为
\[
f_i-m(Y)=f_i-f_{n+1}-\frac{1}{n+1}\sum_{k=1}^n kb_k=\sum_{k=i}^n kb_k.
\]
% [Tutor Note: 上式最右端照原文排印。按 $m(Y)=\frac{1}{n+1}\sum_{j=1}^{n+1}f_j$ 与
% $f_i-f_{i+1}=b_i$ 直接计算，$f_i-m(Y)=\sum_{k=i}^n b_k-\frac{1}{n+1}\sum_{k=1}^n kb_k$，
% 而非 $\sum_{k=i}^n kb_k$；此处原文有笔误（见报告 ERRATA）。要点是 $f_i-m(Y)$
% 只依赖 $b_j$（即只依赖 $\lambda$），这正是它不依赖 $Y$ 的原因。

从杨图
\[
Y=[f_1,f_2,\dots,f_{n+1}]
\]
中取走部分 $[\underbrace{f_{n+1},f_{n+1},\dots,f_{n+1}}_{n+1}]$，我们得到杨图
$Y^{red}=[f_1-f_{n+1},f_2-f_{n+1},\dots,f_n-f_{n+1}]\in\Gamma_{n,\ell}$。
杨图 $Y$ 与 $Y^{red}$ 给出 $\mathfrak{sl}(n+1,\C)$ 的同一个表示。

让我们定义 $\bar\varepsilon_j\in\mathfrak{h}^*$，使
\[
\lambda(Y)+\bar\varepsilon_j=\lambda(Y+\bar\varepsilon_j),
\]
其中对杨图 $Y=[f_1,f_2,\dots,f_{n+1}]$，我们定义
\[
Y+\bar\varepsilon_j=[f_1,\dots,f_{j-1},f_j+1,f_{j+1},\dots,f_{n+1}].
\]
% [Tutor Note: $Y+\bar\varepsilon_j$ 是把第 $j$ 行增加一格所得的杨图；于是
% $\lambda(Y+\bar\varepsilon_j)-\lambda(Y)$ 就是加一格所对应的权，记为
% $\bar\varepsilon_j\in\mathfrak{h}^*$。由 (6.34) 它与 $\varepsilon_j$ 只差
% 一个与 $n+1$ 有关的常数平移，故不必与 $\varepsilon_j$ 混淆。

以下对杨图 $Y$，用 $V_Y$ 记左 $\mathfrak{g}$-模 $V_{\lambda(Y)}$。

现在考虑四点情形 $\mathbf{\Lambda}=(\lambda_4,\lambda(\square),\lambda(\square),\lambda_1)$，
其中 $V_{\mathfrak{g}}^*(\mathbf{\Lambda})\neq0$。

\figomit{（无编号）}{四点共形块的符号图：一条水平线段（表示 $\PP^1$）上有两个
带向下箭头的插入点，各自上方标 $\lambda(\square)$；水平线左端标 $\lambda_4^\dagger$、
中部标 $\mu$、右端标 $\lambda_1$，线上的箭头指向左方。}

定义
\[
I(\mathbf{\Lambda})=\{\mu\in P_+\mid\mathcal V(\mathbf{\Lambda})_\mu^{12}\neq0\}
\supset I_\ell(\mathbf{\Lambda})
=\{\mu\in P_\ell\mid\mathcal V(\mathbf{\Lambda})_\mu^{12}\neq0\},
\]
其中 $\mathcal V(\mathbf{\Lambda})_\mu^{12}
=\mathcal V(\mathbf{v}_2(\mu))\otimes\mathcal V(\mathbf{v}_1(\mu))$。注意对顶点
$\mathbf{v}$，向量空间 $\mathcal V(\mathbf{v})$ 的定义见定义 6.14。在我们的情形，
由 $\mathfrak{sl}(n+1,\C)$ 的表示论有

\begin{proposition}[6.21]
对
\[
\mathbf{v}=\binom{\lambda(\square)}{\mu_2\ \mu_1},
\]
其中 $\mu_1,\mu_2\in P_\ell$，我们有
\[
\mathcal V(\mathbf{v})
=\operatorname{Hom}_{\mathfrak{g}}(V_{\square}\otimes V_{\mu_1},V_{\mu_2}).
\]
\end{proposition}

于是得到
\[
V_{\mathfrak{g}}^*(\mathbf{\Lambda})_\mu^{12}
=\mathcal V(\mathbf{\Lambda})_\mu^{12}.
\]
% [Tutor Note: 命题 6.21 是 $\mathfrak{sl}(n+1,\C)$ 的 Clebsch–Gordan 型陈述：一个
% "block" 型的顶点的空间等于由一个方块（基本表示 $V_\square$）与 $V_{\mu_1}$
% 张成、映到 $V_{\mu_2}$ 的 $\mathfrak{g}$-同态空间。原文只用一句"由表示论有"
% 就带过，跳过的细节见报告 DIFFICULT。

现在把 $\lambda_1$ 表为
\[
\lambda_1=\sum_{j=1}^n b_j\omega_j=\lambda(Y),
\]
其中 $Y=[f_1,f_2,\dots,f_{n+1}]$。

% [原文 p.154]
于是不难证明
\[
\dim V_{\mathfrak{g}}^*(\mathbf{\Lambda})\le2.
\]
下述结果是 $\mathfrak{sl}(n+1,\C)$ 的表示论的一个简单推论。

\begin{lemma}[6.22]
$\dim V_{\mathfrak{g}}^*(\mathbf{\Lambda})=2$ 当且仅当
\begin{equation}
\lambda_4^\dagger=\lambda_1+\bar\varepsilon_i+\bar\varepsilon_j,\quad
b_{i-1}\ge1,\quad b_{j-1}\ge1,\quad i<j.
\tag{6.37}
\end{equation}
而且在此情形下
\[
I(\mathbf{\Lambda})=\{\mu_+=\lambda_1+\bar\varepsilon_i,\
\mu_-=\lambda_1+\bar\varepsilon_j\in P_+\}.
\]
又，$I(\mathbf{\Lambda})=I_\ell(\mathbf{\Lambda})$ 当且仅当
\[
\text{Case }(D2)\quad\epsilon_0(\mathbf{\Lambda})<1,\quad j-i<n,\quad
\text{or}\quad j-i=n,\quad(\theta,\lambda_1)<\ell,
\]
其中
\[
d(\mathbf{\Lambda})=j-i+f_i-f_j,\quad
\epsilon_0(\mathbf{\Lambda})=\frac{d\Lambda)+1}{\kappa},\quad
\kappa=n+1+\ell.
\]
\end{lemma}
% [Tutor Note: 关于引理 6.22 有三点说明。(1) 上式分子原文印作 "$d\Lambda)+1$"，
% 缺一个左括号，应为 $d(\mathbf{\Lambda})+1$，照原文排印（见报告 ERRATA）；
% $\kappa=n+1+\ell$ 是本章反复出现的规范化常数，与 §5.6 的 $2\ell+2g^*$ 是同一类平移。
% (2) 这里以及后面的 $f_i,f_j,m$ 都指上一段那个杨图 $Y$ 的行长与 $m(Y)$：上一页证明了
% $f_i-m(Y)$ 不依赖 $Y$ 的选取（对固定的 $\lambda_1$），所以下面的
% $\gamma_\pm^{(i)}$ 是良定的；原文没有重复这一点，读者容易误以为 $f_i,m$ 是自由参数。
% (3) $I(\mathbf{\Lambda})$ 是全部支配整权 $\mu$ 中的指标集，$I_\ell(\mathbf{\Lambda})$
% 是其中落在标号集 $P_\ell$ 里的部分；条件 $I=I_\ell$ 的意思是没有"权太大"的中间顶点，
% 这正是情形 (D2) 与非 (D2) 的唯一差别，由此立刻得到下面
% $\dim\mathcal V(\mathbf{\Lambda})$ 的取值。

因此我们断言：若 $\mathcal V(\mathbf{\Lambda})\neq0$，则有
\[
\dim\mathcal V(\mathbf{\Lambda})=
\begin{cases}
2 & \text{Case }(D2)\\[2pt]
1 & \text{otherwise}.
\end{cases}
\]

现在考虑情形 (D2)。选取 $V_{\mathfrak{g}}^*(\mathbf{\Lambda})$ 的三组基
$\{U_\pm^{(0)}\}$、$\{U_\pm^{(1)}\}$、$\{U_\pm^{(\infty)}\}$，使它们分别对角化算子
$\Omega_{12}$、$\Omega_{23}$、$\Omega_{13}$。也就是说，
\[
\Omega_{12}U_\pm^{(0)}=\kappa(\gamma_\pm^{(0)}-\Delta_4(\mathbf{\Lambda}))U_\pm^{(0)},
\quad
\Omega_{23}U_\pm^{(1)}=\kappa\gamma_\pm^{(1)}\Omega_{23}U_\pm^{(1)},
\]
\[
\Omega_{13}U_\pm^{(\infty)}=\kappa\gamma_\pm^{(\infty)}\Omega_{13}U_\pm^{(\infty)},
\]
其中
\[
\begin{gathered}
\gamma_+^{(0)}=\frac{1}{\kappa}\Bigl(j-f_j+m-\frac{n}{n+1}\Bigr),\quad
\gamma_+^{(1)}=\frac{n}{\kappa(n+1)},\quad
\gamma_+^{(\infty)}=\frac{1}{\kappa}(i-f_i-m+1),\\
\gamma_-^{(0)}=\frac{1}{\kappa}\Bigl(i-f_i+m-\frac{n}{n+1}\Bigr),\quad
\gamma_-^{(1)}=\frac{-(n+2)}{\kappa(n+1)},\quad
\gamma_-^{(\infty)}=\frac{1}{\kappa}(j-f_j-m+1).
\end{gathered}
\]
% [Tutor Note: 这两行特征值等式的右端原文把算子 $\Omega_{23}$（相应地 $\Omega_{13}$）
% 又写了一遍：$\Omega_{23}U_\pm^{(1)}=\kappa\gamma_\pm^{(1)}\Omega_{23}U_\pm^{(1)}$。
% 按上一行 $\Omega_{12}U_\pm^{(0)}=\kappa(\gamma_\pm^{(0)}-\Delta_4)U_\pm^{(0)}$ 的格式，
% 右端应是 $\kappa\gamma_\pm^{(1)}U_\pm^{(1)}$（见报告 ERRATA）。本译仍照原文排印，
% "对角化"的意思只在去掉重复的算子后才成立。

通过适当地选取这三组基 $U_\pm^{(i)}$，Tsuchiya–Kanie [TK] 与 Kanie [Kan] 证明了
以下定理。

\begin{theorem}[6.23]
KZ 方程 $KZ_1(\mathbf{\Lambda})$ 在 $\zeta=i$（$i=0,1$ 与 $\infty$）处正规化的
基本解 $\Psi_\pm^{(i)}(\zeta)$ 由
\[
(\Psi_+^{(i)}(\zeta),\Psi_-^{(i)}(\zeta))
=(U_+^{(i)},U_-^{(i)})
\begin{pmatrix}
\varphi_{++}^{(i)}(\zeta) & \varphi_{+-}^{(i)}(\zeta)\\
\varphi_{-+}^{(i)}(\zeta) & \varphi_{--}^{(i)}(\zeta)
\end{pmatrix}
\]
给出。
\end{theorem}

% [原文 p.155]
这里，在 $\zeta=0$ 处借助超几何函数（hypergeometric function）
$F(\alpha,\beta,\gamma;\zeta)$，$\varphi_{\bullet\bullet}^{(0)}(\zeta)$ 写为
\[
\begin{aligned}
\varphi_{++}^{(0)}(\zeta)&=\zeta^{\gamma_+^{(0)}}
(1-\zeta)^{\gamma_+^{(1)}}F(\alpha,\beta,\gamma^{(0)};\zeta),\\
\varphi_{+-}^{(0)}(\zeta)&=c_+^{(0)}\zeta^{1+\gamma_+^{(0)}}
(1-\zeta)^{\gamma_+^{(1)}}F(\alpha+1,\beta+1,\gamma^{(0)}+2;\zeta),\\
\varphi_{-+}^{(0)}(\zeta)&=c_-^{(0)}\zeta^{1+\gamma_-^{(0)}}
(1-\zeta)^{\gamma_-^{(1)}}F(-\alpha+1,-\beta+1,-\gamma^{(0)}+2;\zeta),\\
\varphi_{--}^{(0)}(\zeta)&=\zeta^{\gamma_-^{(0)}}
(1-\zeta)^{\gamma_-^{(1)}}F(-\alpha,-\beta,-\gamma^{(0)};\zeta),
\end{aligned}
\]
其中
\[
\alpha=\frac{d(\mathbf{\Lambda})+1}{\kappa},\quad
\beta=\frac{1}{\kappa},\quad
c_\pm^{(0)}=-\frac{\sqrt{d(\mathbf{\Lambda})+1}\sqrt{d(\mathbf{\Lambda})-1}}
{d(\mathbf{\Lambda})(\kappa\pm d(\mathbf{\Lambda}))}.
\]
类似地，在 $\zeta=\infty$ 处我们有
\[
\begin{aligned}
\varphi_{++}^{(\infty)}(\zeta)&=\zeta^{-\gamma_+^{(\infty)}}
\Bigl(1-\frac{1}{\zeta}\Bigr)^{\gamma_+^{(1)}}
F\Bigl(\alpha,\beta,\gamma^{(\infty)};\frac{1}{\zeta}\Bigr),\\
\varphi_{+-}^{(\infty)}(\zeta)&=c_+^{(0)}\zeta^{-1-\gamma_+^{(\infty)}}
\Bigl(1-\frac{1}{\zeta}\Bigr)^{\gamma_+^{(1)}}
F\Bigl(\alpha+1,\beta+1,\gamma^{(\infty)}+2;\frac{1}{\zeta}\Bigr),\\
\varphi_{-+}^{(\infty)}(\zeta)&=c_-^{(0)}\zeta^{-1-\gamma_-^{(\infty)}}
\Bigl(1-\frac{1}{\zeta}\Bigr)^{\gamma_-^{(1)}}
F\Bigl(-\alpha+1,-\beta+1,-\gamma^{(\infty)}+2;\frac{1}{\zeta}\Bigr),\\
\varphi_{--}^{(\infty)}(\zeta)&=\zeta^{-\gamma_-^{(\infty)}}
\Bigl(1-\frac{1}{\zeta}\Bigr)^{\gamma_-^{(1)}}
F\Bigl(-\alpha,-\beta,-\gamma^{(\infty)};\frac{1}{\zeta}\Bigr).
\end{aligned}
\]
% [Tutor Note: 这八行里的 $\gamma^{(0)}$、$\gamma^{(\infty)}$ 原文都省略了下标：按本页
% 上一页的特征值记号，$\varphi_{\pm\bullet}$ 各自的第三参数应当是 $\gamma_\pm^{(0)}$
% 或 $\gamma_\pm^{(\infty)}$（$\bullet$ 指同一行的 $\pm$）。尤其 $\varphi_{-+}$ 行里
% $F$ 的第三参数写成 $-\gamma^{(0)}+2$，负号来自 $\pm=-$ 那一支。照原文排印。
% [Tutor Note: 超几何函数的解在此是 $\zeta^{\gamma}(1-\zeta)^{\gamma^{(1)}}$ 型的
% 两支解（$F$ 的第三参数相差符号），这正是 KZ 方程在正则奇点处两个基本解的
% 标准形式；四个分量 $\varphi_{\bullet\bullet}$ 只是把 $\zeta=0$ 与 $\zeta=\infty$
% 两处的解配对。原文没有给 $F$ 的级数定义，见本章附录。

关于超几何函数见本章附录。由这一定理我们推出情形 (D2) 中的连接矩阵。

\begin{corollary}[6.24]
连接矩阵 $K(\mathbf{\Lambda})$ 由
\[
(\Psi_+^{(0)},\Psi_-^{(0)})
=(\Psi_+^{(\infty)},\Psi_-^{(\infty)})
\begin{pmatrix}
K_+^+ & K_-^+\\
K_+^- & K_-^-
\end{pmatrix}
\]
给出，其中
\[
\begin{aligned}
K_+^+&=-q^{\frac{-1}{2}\bigl(d+\frac{1}{n+1}\bigr)}
\frac{\gamma}{\beta}\cdot\frac{\Gamma(\gamma)\Gamma(-\gamma)}
{\Gamma(\beta)\Gamma(-\beta)},&
K_-^+&=-q^{\frac{-1}{2(n+1)}}\Bigl(\frac{\gamma^2}{\alpha\epsilon}\Bigr)^{1/2}
\frac{\Gamma(-\gamma)^2}{\Gamma(-\alpha)\Gamma(-\epsilon)},\\
K_+^-&=q^{\frac{-1}{2(n+1)}}\Bigl(\frac{\gamma^2}{\alpha\epsilon}\Bigr)^{1/2}
\frac{\Gamma(\gamma)^2}{\Gamma(\alpha)\Gamma(\epsilon)},&
K_-^-&=-q^{\frac{1}{2}\bigl(d+\frac{1}{n+1}\bigr)}
\frac{\gamma}{\beta}\cdot\frac{\Gamma(\gamma)\Gamma(-\gamma)}
{\Gamma(\beta)\Gamma(-\beta)},
\end{aligned}
\]
这里对 $d=d(\mathbf{\Lambda})$ 我们定义
\[
\alpha=\frac{d+1}{\kappa},\quad\beta=\frac{1}{\kappa},\quad
\gamma=\frac{d}{\kappa},\quad\epsilon=\frac{d-1}{\kappa},\quad
q=e^{2\pi\sqrt{-1}/\kappa}.
\]
\end{corollary}
% [Tutor Note: 由上一条定义 $\beta=1/\kappa$、$\gamma=d/\kappa$，故 $\gamma/\beta=d$，
% 于是 $q$ 的指数可写成 $\frac{1}{2}(d+\frac{1}{n+1})=\frac{1}{2}(\gamma/\beta+\frac{1}{n+1})$；
% 连接矩阵的四个元素因此都是 $d$（即 $\Lambda$ 的"距离"）的函数。注意此处的
% $\gamma$ 与前面 $\gamma_\pm^{(i)}$（KZ 特征值）是两套记号，原文重用了 $\gamma$。

我们可以把结果改写成下述形式。

% [原文 p.156]
\begin{theorem}[6.25]
对满足 $(D2)$ 的 $\mathbf{\Lambda}=(\lambda_4,\lambda(\square),\lambda(\square),\lambda_1)$，
连接矩阵 $K(\mathbf{\Lambda})$ 可写成下述形式。
\[
K(\mathbf{\Lambda})=q^{\frac{-(n+2)}{2(n+1)}}
\begin{pmatrix}
\gamma_+^{-1} & \\[2pt]
 & \gamma_-^{-1}
\end{pmatrix}
\begin{pmatrix}
\displaystyle\frac{-1}{[d]} &
\displaystyle\frac{\sqrt{q[d+1][d-1]}}{[d]}\\[10pt]
\displaystyle\frac{\sqrt{q[d+1][d-1]}}{[d]} &
\displaystyle\frac{q^d}{[d]}
\end{pmatrix}
\begin{pmatrix}
\gamma_+ & \\[2pt]
 & \gamma_-
\end{pmatrix}
\]
\end{theorem}
% [Tutor Note: 定理 6.25 里左右两个对角矩阵原文把两个元素并排印出、没有分隔，
% 本译按对角矩阵排出（空位留白）。中间的对称矩阵是 $\mathrm{SL}(2,\Z)$ 或
% 量子群 $U_q(\mathfrak{sl}_2)$ 的"量子 $6j$-符号"型矩阵；$[d]$、$[d\pm1]$ 是
% 下面的 $q$-整数，当 $q=1$ 时它们退化为通常的整数，整个 $K$ 退化为
% 相应的经典连接系数。

其中
\[
d=d(\mathbf{\Lambda}),\quad
\gamma_\pm=\frac{\Gamma(\pm\gamma)}
{\sqrt{\Gamma(\pm\alpha)\Gamma(\pm\epsilon)}}
\]
以及
\[
[j]=
\begin{cases}
\displaystyle\frac{q^j-1}{q-1}, & q\neq1,\\[10pt]
j, & q=1.
\end{cases}
\]
% [Tutor Note: $\gamma_\pm$ 的定义里用了"$\pm$"的整体缩写：上行取 $\gamma_+$ 与
% $\Gamma(+\gamma),\Gamma(+\alpha),\Gamma(+\epsilon)$，下行取 $\gamma_-$ 与
% 相应带负号的 Gamma 值。$\alpha,\epsilon$ 已在本页前面定义为 $(d\pm1)/\kappa$，
% 所以 $\gamma_\pm$ 只是 $d$ 与 $\kappa$ 的函数（再看 $\kappa=n+1+\ell$）。

对其它情形，带坐标以及 $\mathbf{\Lambda}=(\lambda_4,\lambda(\square),
\lambda(\square),\lambda_1)$ 的四点 $\PP^1$ 的共形块维数为一，而连接矩阵则
由 Tsuchiya–Kanie [TK] 与 Kanie [Kan] 给出下述结果。
% [Tutor Note: 原文"with coordinates and $\Lambda=(\dots)$"之间缺了坐标的列举
% （按上下文应为 $0,1,\infty,\zeta$ 之类），照原文排印（见报告 ERRATA）。

\begin{theorem}[6.26]
若 $\mathbf{\Lambda}=(\lambda_4,\lambda(\square),\lambda(\square),\lambda_1)$
不属于情形 $(D2)$，则连接矩阵 $K(\mathbf{\Lambda})$ 或为
\[
K(\mathbf{\Lambda})=-q^{\frac{-(n+2)}{2(n+2)}}
\]
或为
\[
K(\mathbf{\Lambda})=-q\cdot q^{\frac{-(n+2)}{2(n+2)}}.
\]
\end{theorem}
% [Tutor Note: 定理 6.26 的指数分母原文印作 $2(n+2)$，与定理 6.25 的 $2(n+1)$
% 不一致，应为 $2(n+1)$（见报告 ERRATA）；两式相差一个因子 $q$，即把
% 基本解的规范化再乘一个常数，故"要么…要么…"只是规范化选取的自由度。

\noindent\textbf{6.5.2 单值性表示.}\quad
让我们考虑 $(N+2)$ 点 $\PP^1$ 的族的共形块丛
\[
(\PP^1;\infty,z_N,z_{N-1},\dots,z_1,0;
1/\zeta,\zeta-z_N,\dots,\zeta-z_1,\zeta),
\]
它定义在 $X_N$ 上，其参数为
\begin{equation}
\mathbf{\Lambda}_\tau=(\tau,\lambda(\square),\dots\lambda(\square),0).
\tag{6.38}
\end{equation}
如 §6.3 所示，一个 $N$ 点函数给出该丛的一个多值截面。

让我们考虑多值截面空间上的单值性表示。首先注意，以
$(N,N-1,\dots,1)$ 为基点的 $X_N$ 的基本群是 $N$ 股限制辫群（restricted braid
group）。设 $\overline{X_N}$ 为 $X_N$ 关于第 $N$ 个对称群的商流形
$X_N/S_N$。$\overline{X_N}$ 的基本群是 $N$ 股辫群（braid group）$B_N$。设
$\sigma_i(t)$，$1\le i\le N-1$，为 $\overline{X_N}$ 中由
\[
\sigma_i(t)=\Bigl(N,N-1,\dots,i+\frac{1}{2}(1+e^{\pi\sqrt{-1}t},\dots,1\Bigr),
\quad t\in[0,1]
\]
定义的道路。
% [Tutor Note: 上式内层括号原文只开不闭：$i+\frac{1}{2}(1+e^{\pi\sqrt{-1}t}$ 之后
% 直接接 $\dots,1$，缺一个右括号，应为 $i+\frac{1}{2}(1+e^{\pi\sqrt{-1}t})$。
% 其含义是第 $i$ 与第 $i+1$ 个坐标绕彼此半圈互换，$t$ 从 $0$ 到 $1$ 时
% $e^{\pi\sqrt{-1}t}$ 从 $1$ 走到 $-1$，两个坐标恰好交换。

于是集合 $\{\sigma_i(t),i=1,\dots,N-1\}$ 是辫群 $B_N$ 的一组生成元，它们满足
基本关系：
\[
\sigma_i\sigma_j=\sigma_j\sigma_i\quad\text{若 }|i-j|\ge2,\qquad
\sigma_i\sigma_{i+1}\sigma_i=\sigma_{i+1}\sigma_i\sigma_{i+1},\quad
1\le i\le N-2.
\]
% [Tutor Note: 这里 $X_N$ 是 $N$ 个互异点的有序构形空间，故其基本群是纯辫群；
% 取商 $X_N/S_N$ 后才是通常的辫群 $B_N$。原文用"restricted braid group"
% 指前者（限定在原点附近取基点的纯辫群），与后文的 $B_N$ 不是同一个群，
% 单值性表示正是在 $B_N$ 上取的。

% [原文 p.157]
让我们考虑 $V_{\mathfrak{g}}^*(\mathbf{\Lambda}_\tau)$-值函数的 KZ 方程
$KZ(N;\tau)$ 与 $\ell C$-约束 $\ell C(N;\tau)$，其中标号
$\mathbf{\Lambda}_\tau$ 由 (6.38) 给出。
\[
KZ(N:\tau):\quad
\Bigl(\kappa\frac{\partial}{\partial z_i}
-\sum_{k=1,k\neq i}^{N}\frac{\Omega_{ik}}{z_i-z_k}\Bigr)
\Psi(z_N,\dots,z_1)=0,\quad i=1,\dots,N
\]
\begin{align*}
\ell C(N;\tau):\quad
&\sum_{\vec m_i}\binom{L_i}{\vec m_i}
\prod_{k=1,k\neq i}^{N}(z_k-z_i)^{-m_k}\Psi(z_n,\dots,z_1)\\
&\times(u_{N+1},X_\theta^{m_N}u_N,\dots,|\lambda(\square)\rangle,
\dots,X_\theta^{m_1}u_1)=0,\quad i=1,\dots,N\\
&\sum_{\vec m_{N+1}}\binom{L_{N+1}}{\vec m_{N+1}}
\Psi(z_n,\dots,z_1)(|\mu^\dagger\rangle,X_\theta^{m_N}u_N,
\dots,X_\theta^{m_1}u_1)=0,
\end{align*}
即对任意 $u_N\in V_{\tau^\dagger}$ 与 $u_i\in V_{\square}$（$i=1,\dots,N$）成立，
其中
\[
\vec m_i=(m_N,\dots,\widehat{m_i},\dots,m_1)\in(\Z_{\ge0})^{N-1},\quad
\vec m_{N+1}=(m_n,\dots,m_1)\in(\Z_{\ge0})^{N},
\]
且
\[
|\vec m_i|=L_i=\ell-(\theta,\lambda(\square))+1,\quad i=1,\dots,N,\qquad
|\vec m_{N+1}|=L_{N+1}=\ell-(\theta,\tau^\dagger)+1.
\]
% [Tutor Note: $\ell C$-约束的三个式子原文有几处小写 $n$（应是 $N$），以及
% 第一个式子里 $u_{N+1}$ 与正文"$u_N\in V_{\tau^\dagger}$"的指标不一致；
% 末式左端的 $|\mu^\dagger\rangle$ 与第二式里的 $|\lambda(\square)\rangle$ 相对，
% 应理解为另一端（$\tau$ 的共轭）的插入。以上均照原文排印，见报告 ERRATA。
% $\widehat{m_i}$ 表示删去第 $i$ 个分量，故 $\vec m_i$ 有 $N-1$ 个分量、
% $\vec m_{N+1}$ 有 $N$ 个分量；$|\vec m|=\sum_j m_j$。

于是 $\sigma_i$ 的单值性表示 $\pi_{N,\tau}$ 由联合系统 $KZ(N;\tau)$ 与
$\ell C(N;\tau)$ 的解空间 $W(N;\tau)$ 沿道路
\[
\begin{aligned}
\gamma_i(t)&=\Bigl(z_N,z_{N-1},\dots,
\frac{z_{i+1}+z_i}{2}+\frac{z_{i+1}-z_i}{2}\,e^{\pi\sqrt{-1}t},\\
&\qquad
\frac{z_{i+1}+z_i}{2}-\frac{z_{i+1}-z_i}{2}\,e^{\pi\sqrt{-1}t},
\dots,z_1\Bigr),\quad t\in[0,1]
\end{aligned}
\]
的解析延拓给出。

§6.4 的论证表明解空间 $W(N;\tau)$ 同构于
\[
\mathcal V(N;\tau)
=\sum_{\vec\mu\in(P_\ell)^{N-1}}\mathcal V(N;\tau)_{\vec\mu},
\quad \vec\mu=(\mu_{N-1},\dots,\mu_1)\in(P_\ell)^{N-1},
\]
其中
\[
\mathcal V(N;\tau)_{\vec\mu}
=\mathcal V\Bigl(\binom{\lambda(\square)}{\tau\ \mu_{N-1}}\Bigr)
\otimes\cdots\otimes
\mathcal V\Bigl(\binom{\lambda(\square)}{\mu_i\ \mu_{i-1}}\Bigr)
\otimes\cdots\otimes
\mathcal V\Bigl(\binom{\lambda(\square)}{\mu_1\ 0}\Bigr).
\]

由命题 6.20，$\pi_{N,\tau}(\sigma_i)$ 可用定理 6.25 与定理 6.26 算出。

注意 $\operatorname{Hom}_{\mathfrak{g}}(V_{\square}\otimes V_{\mu_j},
V_{\mu_{j+1}})$ 是零维或一维的。向量空间
$\operatorname{Hom}_{\mathfrak{g}}(V_{\square}\otimes V_{\mu_j},
V_{\mu_{j+1}})$ 是一维的，当且仅当对应 $\mu_j$ 的杨图
$Y_{\mu_j}\in\Gamma_{n,\ell}$（带有 $j$ 个格子）可由对应 $\mu_{j+1}$ 的杨图
$Y_{\mu_{j+1}}$ 删去一个格子得到。于是若置
\[
\begin{aligned}
\mathcal P_\ell(N;\tau)=\{\vec{\mathbf p}
&=(\lambda_N,\lambda_{N-1},\dots,\lambda_1,\lambda_0)\mid
\lambda_N=\tau,\ \lambda_0=0,\\
&\quad\lambda_i\in P_\ell,\ 1\le i\le N,\
\lambda_i=\lambda_{i-1}+\bar\epsilon_j\text{ 对某个 }j\},
\end{aligned}
\]
则由定理 6.18，解空间 $W(N;\tau)$ 有一组基
\[
\{\Psi_{\vec p}(z_n,\dots,z_1)\mid\vec p\in\mathcal P_\ell(N;\tau)\},
\]
其中
\[
\Psi_{\vec p}(z_n,\dots,z_1)=\langle\Phi_{\mathbf{v}_N}\cdots\Phi_{\mathbf{v}_1}\rangle.
\]

% >>>END

